% Auto-generated from book/HOW_TO_USE.md — do not edit by hand
% Regenerated by scripts/md_to_latex.py

\chapter*{How to Use the Figures and Code}
\label{ch:how_to_use}
\addcontentsline{toc}{chapter}{How to Use the Figures and Code}
\markboth{How to Use the Figures and Code}{}

This short front-matter page is the practical companion to the Preface. It fixes the conventions used in every chapter so figures render correctly and labs stay reproducible. For navigation, start with root \texttt{TOC.\allowbreak{}md}; for Hatcher's book in parallel, use \texttt{HATCHER\_\allowbreak{}MAP.\allowbreak{}md}.

\bigskip
\noindent\rule{\textwidth}{0.4pt}
\bigskip

\section*{Claim labels}
\label{ch:how_to_use:claim-labels}

Every major claim in the manuscript is one of four kinds:

\begin{center}
\small
\setlength{\tabcolsep}{4pt}
\begin{tabularx}{\textwidth}{@{}>{\raggedright\arraybackslash}p{0.12\textwidth}>{\raggedright\arraybackslash}X>{\raggedright\arraybackslash}X@{}}
\toprule
Label & Meaning & Example \\
\midrule
\textbf{Theorem} & Proved in this book or classical & Four-square theorem; Hopf linking; Hurwitz class number 1 \\
\textbf{Model} & Consistent mathematical or physical construction; not claimed as nature's unique law & Gauged Hopf lattice adjacency; flux topographs; \(Z\mapsto\) proxy \\
\textbf{Hypothesis} & Observational claim needing validation or falsification & H1a--H1e: \(W_g=350/\pi\) per domain (not one bundled clock) \\
\textbf{Software fact} & True of the current code or portal implementation & \texttt{map\_\allowbreak{}z\_\allowbreak{}to\_\allowbreak{}flywheel} returns \texttt{stability\_\allowbreak{}score}; there is \textbf{no} \texttt{magic\_\allowbreak{}flag} key \\
\bottomrule
\end{tabularx}
\end{center}


\textbf{Rule of thumb.} Geometry and classical algebra \(\rightarrow\) Theorem. Lattice / topograph design choices \(\rightarrow\) Model. Portal numerics about the physical world \(\rightarrow\) Hypothesis. ``What this function returns today'' \(\rightarrow\) Software fact.

Validation protocols for Hypotheses live in Chapter 10 (Table \textbf{T4}) and \texttt{lib/\allowbreak{}validation.\allowbreak{}py}.

\bigskip
\noindent\rule{\textwidth}{0.4pt}
\bigskip

\section*{Repository layout}
\label{ch:how_to_use:repository-layout}

\begin{center}
\small
\setlength{\tabcolsep}{4pt}
\begin{tabularx}{\textwidth}{@{}>{\raggedright\arraybackslash}X>{\raggedright\arraybackslash}X@{}}
\toprule
Path & Role \\
\midrule
\texttt{book/\allowbreak{}*.\allowbreak{}md} & Manuscript (Preface, HOW\_TO\_USE, Chapters 0--10) \\
\texttt{book/\allowbreak{}figures/\allowbreak{}} & Static figures \(\cdot\) \texttt{ATTRIBUTION.\allowbreak{}md} \\
\texttt{lib/\allowbreak{}} & Pedagogical helpers (not a replacement for Kingdom Come) \\
\texttt{scripts/\allowbreak{}generate\_\allowbreak{}chN\_\allowbreak{}figures.\allowbreak{}py} & Regenerate chapter figures \\
\texttt{TOC.\allowbreak{}md} \(\cdot\) \texttt{HATCHER\_\allowbreak{}MAP.\allowbreak{}md} \(\cdot\) \texttt{notes/\allowbreak{}open\_\allowbreak{}problems.\allowbreak{}md} & Navigation and open research edges \\
\url{https://github.com/kinaar8340/qga} & Canonical remote for this book \\
\bottomrule
\end{tabularx}
\end{center}


\textbf{Live portal and domain code} (separate repo):

\begin{center}
\small
\setlength{\tabcolsep}{4pt}
\begin{tabularx}{\textwidth}{@{}>{\raggedright\arraybackslash}X>{\raggedright\arraybackslash}X@{}}
\toprule
Resource & Location \\
\midrule
Kingdom Come & \url{https://github.com/kinaar8340/kingdom_come} \(\cdot\) extra \texttt{.\allowbreak{}[portal]} \\
Shared Hopf / quaternion core & \texttt{flux-hopf-lib} (same extra) \\
\bottomrule
\end{tabularx}
\end{center}


\bigskip
\noindent\rule{\textwidth}{0.4pt}
\bigskip

\section*{Running the labs}
\label{ch:how_to_use:running-the-labs}

\subsection*{One install path}
\label{ch:how_to_use:one-install-path}

From the clone root (any path; do not hard-code \texttt{\textasciitilde{}\allowbreak{}/\allowbreak{}Projects/\allowbreak{}.\allowbreak{}.\allowbreak{}.\allowbreak{}}):

\begin{lstlisting}[style=qga]
python3 -m pip install -e .
python3 -c "from lib.hopf_lattice import HURWITZ_UNITS, hopf_map; print(len(HURWITZ_UNITS))"
# -> 24
\end{lstlisting}

Labs that import \texttt{kingdom} or \texttt{flux\_\allowbreak{}hopf\_\allowbreak{}lib} need the documented extra (PyPI names \texttt{kingdom-come} / \texttt{flux-hopf-lib}):

\begin{lstlisting}[style=qga]
python3 -m pip install -e ".[portal]"
\end{lstlisting}

Then:

\begin{lstlisting}[style=qga,language=python]
from lib.hopf_lattice import HURWITZ_UNITS, hopf_map, sample_angle_lattice
from lib.flux_topograph import build_flux_topograph
from lib.composition import compose_flywheels
from lib.quaternion_algebra import HurwitzOrder, left_ideal_class_set
from lib.validation import table_t4_checklist, run_table_t4_demo, h2_ablation_table
\end{lstlisting}

\begin{center}
\small
\setlength{\tabcolsep}{4pt}
\begin{tabularx}{\textwidth}{@{}>{\raggedright\arraybackslash}p{0.12\textwidth}>{\raggedright\arraybackslash}X>{\raggedright\arraybackslash}X@{}}
\toprule
Package & Chapters & Purpose \\
\midrule
\texttt{lib.\allowbreak{}hopf\_\allowbreak{}lattice} & 3--4 & Hurwitz units, angle lattice, adjacency (OP1), gauge sequences \\
\texttt{lib.\allowbreak{}flux\_\allowbreak{}topograph} & 5--7 & Topographs, separators, classification, Magic Island scores \\
\texttt{lib.\allowbreak{}composition} & 8 & Flywheel composition, class-group analogue (OP6) \\
\texttt{lib.\allowbreak{}quaternion\_\allowbreak{}algebra} & 9 & Quaternion algebras, Hurwitz order, toy class number \\
\texttt{lib.\allowbreak{}validation} & 10 & Table T4 checklist, Fisher combine, \(W_g\) diagnostics \\
\bottomrule
\end{tabularx}
\end{center}


\subsection*{Kingdom Come / portal modules}
\label{ch:how_to_use:kingdom-come-portal-modules}

After \texttt{pip install -e ".\allowbreak{}[portal]"}:

\begin{lstlisting}[style=qga]
python3 -c "from kingdom.core.quaternion import Quaternion; print('ok')"
\end{lstlisting}

Do not set \texttt{PYTHONPATH=\textasciitilde{}\allowbreak{}/\allowbreak{}Projects/\allowbreak{}kingdom/\allowbreak{}src:\allowbreak{}.\allowbreak{}.\allowbreak{}.\allowbreak{}} as the official path.

\begin{center}
\small
\setlength{\tabcolsep}{4pt}
\begin{tabularx}{\textwidth}{@{}>{\raggedright\arraybackslash}X>{\raggedright\arraybackslash}X@{}}
\toprule
Chapter & Primary modules \\
\midrule
1 & \texttt{kingdom.\allowbreak{}core.\allowbreak{}quaternion} \\
2 & \texttt{kingdom.\allowbreak{}core.\allowbreak{}hopf}, \texttt{kingdom.\allowbreak{}viz.\allowbreak{}hopf\_\allowbreak{}plotly} \\
3--4 & \texttt{lib.\allowbreak{}hopf\_\allowbreak{}lattice}, \texttt{kingdom.\allowbreak{}simulations.\allowbreak{}lattice.\allowbreak{}TwoGyroLattice} \\
5--6 & \texttt{lib.\allowbreak{}flux\_\allowbreak{}topograph} \\
7 & \texttt{kingdom.\allowbreak{}core.\allowbreak{}flux\_\allowbreak{}flywheel} \(\cdot\) \texttt{lib.\allowbreak{}flux\_\allowbreak{}topograph.\allowbreak{}stability\_\allowbreak{}landscape\_\allowbreak{}z} \\
8 & \texttt{lib.\allowbreak{}composition} \\
9 & \texttt{lib.\allowbreak{}quaternion\_\allowbreak{}algebra} \\
10 & \texttt{lib.\allowbreak{}validation} \\
\bottomrule
\end{tabularx}
\end{center}


\subsection*{API honesty (read once, save time)}
\label{ch:how_to_use:api-honesty-read-once-save-time}

These mismatch the sketch names that appear in early brainstorms; the manuscript and labs follow the \textbf{live} APIs:

\begin{center}
\small
\setlength{\tabcolsep}{4pt}
\begin{tabularx}{\textwidth}{@{}>{\raggedright\arraybackslash}X>{\raggedright\arraybackslash}X@{}}
\toprule
Do use & Do not assume \\
\midrule
\texttt{stability\_\allowbreak{}score}, \texttt{stability\_\allowbreak{}class}, \texttt{is\_\allowbreak{}noble\_\allowbreak{}gas} & \texttt{magic\_\allowbreak{}flag} \\
\texttt{TwoGyroLattice.\allowbreak{}step\_\allowbreak{}frame()} & \texttt{.\allowbreak{}step()} \\
\texttt{hopf\_\allowbreak{}map(x1,x2,x3,x4)}, \texttt{hopf\_\allowbreak{}map\_\allowbreak{}quaternion}, \texttt{Quaternion.\allowbreak{}hopf\_\allowbreak{}image()} & \texttt{hopf\_\allowbreak{}map(q)} as a Quaternion object \\
\texttt{sample\_\allowbreak{}fiber(.\allowbreak{}.\allowbreak{}.\allowbreak{})} \(\rightarrow\) \textbf{dict} of arrays & bare list of quaternions \\
\texttt{map\_\allowbreak{}z\_\allowbreak{}to\_\allowbreak{}flywheel} / \texttt{map\_\allowbreak{}z\_\allowbreak{}to\_\allowbreak{}flywheel\_\allowbreak{}extended} & undocumented extra keys \\
\bottomrule
\end{tabularx}
\end{center}


\bigskip
\noindent\rule{\textwidth}{0.4pt}
\bigskip

\section*{Figures}
\label{ch:how_to_use:figures}

\begin{enumerate}
\item \textbf{Paths.} In chapter markdown, figure paths are relative to \texttt{book/\allowbreak{}}:  
\end{enumerate}

\texttt{figures/\allowbreak{}fig0\_\allowbreak{}1\_\allowbreak{}hopf\_\allowbreak{}linked\_\allowbreak{}fibers.\allowbreak{}png}.

\begin{enumerate}
\item \textbf{Index.} Full figure list: \texttt{book/\allowbreak{}FIGURES.\allowbreak{}md}.  
\item \textbf{Attribution.} Sources and regenerate commands: \texttt{book/\allowbreak{}figures/\allowbreak{}ATTRIBUTION.\allowbreak{}md}.  
\item \textbf{Regenerate.} From repo root run \texttt{python3 scripts/\allowbreak{}generate\_\allowbreak{}ch1\_\allowbreak{}figures.\allowbreak{}py}. Chapters 7 and 10 may need the portal extra: \texttt{python3 scripts/\allowbreak{}generate\_\allowbreak{}ch7\_\allowbreak{}figures.\allowbreak{}py}.
\item \textbf{Origins.} Chapter 0 partly vendors Kingdom Come assets; later chapters are mostly generated by the scripts above.
\end{enumerate}

Naming pattern: \texttt{fig\{N\}\_\allowbreak{}\{k\}\_\allowbreak{}.\allowbreak{}.\allowbreak{}.\allowbreak{}} for main figures, \texttt{aux\{N\}\_\allowbreak{}\{k\}\_\allowbreak{}.\allowbreak{}.\allowbreak{}.\allowbreak{}} for auxiliary figures, \texttt{aux\_\allowbreak{}z\_\allowbreak{}map\_\allowbreak{}*.\allowbreak{}png} for Z-map stills.

\bigskip
\noindent\rule{\textwidth}{0.4pt}
\bigskip

\section*{Notation (minimal sheet)}
\label{ch:how_to_use:notation-minimal-sheet}

\begin{center}
\small
\setlength{\tabcolsep}{4pt}
\begin{tabularx}{\textwidth}{@{}>{\raggedright\arraybackslash}X>{\raggedright\arraybackslash}X@{}}
\toprule
Symbol & Meaning \\
\midrule
\(\mathbb{H}\), \(S^3\) & Quaternions; unit quaternions \\
\(h:S^3\to S^2\) & Hopf map \\
\(\Lambda_0\) & 24 Hurwitz units on \(S^3\) \\
\(\Lambda_{\mathrm{ang}}\) & Angle-sampled lattice (visualization density) \\
\(\Lambda_{\mathrm{dyn}}\) & Dynamical two-gyro sites (portal) \\
\(W_g\) & \(350/\pi\) (not \(350\pi\)) topological clock (\textbf{Hypothesis} H1a--H1e; Ch. 10) \\
OP1--OP6 & Open problems \(\cdot\) \texttt{notes/\allowbreak{}open\_\allowbreak{}problems.\allowbreak{}md} \\
\bottomrule
\end{tabularx}
\end{center}


\bigskip
\noindent\rule{\textwidth}{0.4pt}
\bigskip

\section*{Gradio portal \(\leftrightarrow\) book}
\label{ch:how_to_use:gradio-portal-book}

\begin{center}
\small
\setlength{\tabcolsep}{4pt}
\begin{tabularx}{\textwidth}{@{}>{\raggedright\arraybackslash}X>{\raggedright\arraybackslash}X@{}}
\toprule
Portal tab & Primary chapters \\
\midrule
\textbf{Book Mode} & \textbf{0--10 map} (chapter picker \(\rightarrow\) jump + mini demos) \\
Hopf Visualizer & 0, 2 \\
Lattice Simulator & 3--4 \\
Flux Flywheel & 0, 5--7, 10 \\
Observations & 10 \\
The Model / theory & Preface, 0, 8--9 (narrative; labs offline) \\
\bottomrule
\end{tabularx}
\end{center}


\subsection*{Book Mode (interactive companion)}
\label{ch:how_to_use:book-mode-interactive-companion}

In the Kingdom Come portal, open the \textbf{Book Mode} tab (between Help and Hopf Visualizer):

\begin{enumerate}
\item Choose a manuscript chapter from the dropdown --- summary, claim focus, and linked live tab appear.
\item Press \textbf{Open linked live tab} to jump the main tab bar to Hopf / Lattice / Flux / Observations / Home.
\item Use the Part subtabs for \textbf{mini demos} (Classic Hopf snapshot, short lattice comparison, Flux Flywheel at a chosen \(Z\)) without leaving Book Mode.
\end{enumerate}

\begin{center}
\small
\setlength{\tabcolsep}{4pt}
\begin{tabularx}{\textwidth}{@{}>{\raggedright\arraybackslash}X>{\raggedright\arraybackslash}X@{}}
\toprule
Book Part & Live widgets \\
\midrule
I Foundations (Ch. 1--2) & Hopf Visualizer \\
II Farey lift (Ch. 3--4) & Lattice Simulator \\
III Forms / \(Z\)-map (Ch. 5--7) & Flux Flywheel \\
IV Arithmetic (Ch. 8--9) & Home \(\cdot\) The Model (+ \texttt{lib/\allowbreak{}composition}, \texttt{lib/\allowbreak{}quaternion\_\allowbreak{}algebra}) \\
V Observations (Ch. 10) & Observations \\
\bottomrule
\end{tabularx}
\end{center}


Source of the map: \texttt{kingdom/\allowbreak{}app/\allowbreak{}pages/\allowbreak{}book\_\allowbreak{}mode.\allowbreak{}py}. Portal: \url{https://github.com/kinaar8340/kingdom_come} \(\cdot\) Space: \url{https://huggingface.co/spaces/kinaar111/kingdom}

\bigskip
\noindent\rule{\textwidth}{0.4pt}
\bigskip

\section*{Parallel reading with Hatcher}
\label{ch:how_to_use:parallel-reading-with-hatcher}

Keep \emph{Topology of Numbers} open beside this book. Full map: root \texttt{HATCHER\_\allowbreak{}MAP.\allowbreak{}md}.

\begin{center}
\small
\setlength{\tabcolsep}{4pt}
\begin{tabularx}{\textwidth}{@{}>{\raggedright\arraybackslash}X>{\raggedright\arraybackslash}X@{}}
\toprule
Hatcher TN & This book \\
\midrule
Ch. 0 Preview & Ch. 0 \\
Ch. 1--3 Farey, CF, symmetries & Ch. 3--4 (after Ch. 1--2 foundations) \\
Ch. 4--6 Forms, classification, representations & Ch. 5--7 \\
Ch. 7--8 Class group, quadratic fields & Ch. 8--9 \\
\emph{(no TN counterpart)} & Ch. 1--2 foundations; Ch. 10 observations \\
\bottomrule
\end{tabularx}
\end{center}


\bigskip
\noindent\rule{\textwidth}{0.4pt}
\bigskip

\section*{Appendices (expanded back matter)}
\label{ch:how_to_use:appendices-expanded-back-matter}

Long listings and reference material live in appendices so the main chapters stay readable:

\begin{center}
\small
\setlength{\tabcolsep}{4pt}
\begin{tabularx}{\textwidth}{@{}>{\raggedright\arraybackslash}X>{\raggedright\arraybackslash}X@{}}
\toprule
App & Content \\
\midrule
\textbf{A} & Terminology and notation \\
\textbf{B} & Open Problems OP1--OP6 (full statements) \\
\textbf{C} & Laboratory code reference (full listings) \\
\textbf{D} & Table T4 validation protocols (full) \\
\textbf{E} & Figure atlas \\
\textbf{F} & Hatcher parallel dictionary \\
\bottomrule
\end{tabularx}
\end{center}


Chapters keep short lab ``calls'' and point here for copy-paste scripts.

\bigskip
\noindent\rule{\textwidth}{0.4pt}
\bigskip

\section*{Suggested first hour}
\label{ch:how_to_use:suggested-first-hour}

\begin{enumerate}
\item Read the Preface and this page.  
\item Open Chapter 0; confirm figures render under \texttt{book/\allowbreak{}figures/\allowbreak{}}.  
\item Run \texttt{from lib.\allowbreak{}hopf\_\allowbreak{}lattice import HURWITZ\_\allowbreak{}UNITS} (expect 24).  
\item Optionally open the Gradio Hopf Visualizer and match panels to Figures 0.1--0.3 / Chapter 2.  
\item When ready for research edges, skim \textbf{Appendix B} (or \texttt{notes/\allowbreak{}open\_\allowbreak{}problems.\allowbreak{}md}).
\end{enumerate}

\bigskip
\noindent\rule{\textwidth}{0.4pt}
\bigskip


